\scanpage{49}
\uline{Ainsi, on trouve que, \inserted{\uline{\unsure{posant|suivant}} \unsure{a)},} les complexes} $g^*(K^\bullet)$, \uline{pour les différentes sections} $g$ \uline{de} $M$ \uline{sur} $S$ \uline{passant par} $t_0$ \uline{sont can. isom. entre eux (au sens des cat. dérivées), donc les} $g^*(\mathcal{H}^i)$
(où $\mathcal{H}^i=\mathcal{H}^i(K^\bullet)$\inserted{$=\mathbb{R}^if_*(\underline{\Omega}^\bullet_{Y/M})$})
\uline{sont canoniquement isom. entre eux}, \deleted{\unsure{puis}} \uline{mod torsion}.
La donnée d'un tel système transitif d'isomorphismes peut être considérée comme une \uline{structure remarquable sur} $K^{\unsure{\bullet|*}}=\mathbb{R}f_*(\underline{\Omega}^\bullet_{Y/M})$ [\uline{en le point} $t_0$ relativement à la base $S$], ou sur les $\mathcal{H}^i=\mathbb{R}^if_*(\Omega^\bullet_{Y/M})$,
\deleted{Une telle}
\deleted{structure est \illegible{2} autrement}
d'espèce voisine à une \uline{structure de descente}.
Une structure de \deleted{\unsure{cet}} ce dernier type, \inserted{sur un Module $\underline{E}$ disons,} à supposer qu'elle soit effective, i.e. définisse un isom
\[
\underline{E}\ \deleted{\unsure{\mathcal{H}^i}}\simeq h^*(\underline{E}_0)
\]
($\underline{E}_0$ un Module sur $S$], définit une

\margin{48}
